\chapter{Preferred Leader Election for Mako}

Preferred leader election solves a real mismatch between standard Raft and
Mako's execution model. Standard Raft can elect any replica whose log is safe
to lead. That is correct for Raft safety, but Mako's benchmark and transaction
path are leader-oriented: the throughput-producing process is expected to
submit transaction logs through the leader-side replication path.

If Raft elects a different process, the benchmark no longer measures the
intended system. The local workers may submit into a non-leader, or the result
may depend on accidental timeout timing. That is not a fair Paxos/Raft
comparison and not a clean deployment story.

Preferred leader election makes leadership placement explicit.

\section{Motivation}

In the benchmark configuration, \code{localhost} is the leader-side process. Its
workers execute transactions, build replicated logs, and submit those logs.
Follower processes such as \code{p1} and \code{p2} primarily replay committed logs. This
asymmetry matches the original Paxos setup.

Raft's default election mechanism does not know about this Mako role. It only
knows which servers are alive, which logs are up to date, and which candidate
can win a majority. Preferred leader election adds placement policy around
Raft without changing Raft safety.

The benchmark motivation is immediate, but the same design issue appears in
real deployments. A system may want leaders near clients, near transaction
coordinators, or on machines provisioned for leader-side CPU and network load.
Mako makes placement especially visible because leadership is connected to
transaction execution and log generation.

\section{Design}

The implementation assigns a preferred replica for each Raft group. During
startup, the preferred replica uses a shorter election timeout, while
non-preferred replicas use longer startup/failover timeouts. This biases the
common startup case toward the desired leader.

If the preferred replica is unavailable, normal Raft failover still works.
Another replica can time out, become a candidate, collect votes, and lead. The
system does not sacrifice availability to preserve placement.

When the preferred replica returns, leadership can transfer back only after
the preferred replica has caught up enough to lead safely. Preference is not
authority. The preferred replica still has to satisfy Raft's normal election
and log-safety conditions.

\begin{figure}[htbp]
\centering
\includegraphics[width=0.96\linewidth]{figures/manual/fig03_preferred_leader_election.pdf}
\caption[Preferred leader election]{Preferred leader election aligns Raft leadership with Mako's leader-oriented submission path while preserving ordinary Raft failover.}
\label{fig:preferred-leader-election}
\end{figure}

\section{Safety Boundary}

Preferred leader election changes placement policy, not consensus safety. A
server still grants at most one vote per term. Candidates still satisfy Raft's
log up-to-date checks. Committed entries still require the normal majority
rule. A preferred replica cannot lead unless it wins a valid Raft election.

The additional engineering condition is safe catch-up before failback. A
non-preferred leader should not step aside merely because the preferred
replica is alive. It should transfer only when the preferred replica has the
committed prefix needed to lead safely. If the preferred replica is behind,
the current leader continues to serve.

This boundary is central to the contribution. Preferred leader election biases
leadership toward a Mako-compatible process. It does not guarantee that the
preferred replica is always leader, and it does not weaken Raft's safety
rules.

\section{Validation}

Preferred leader validation is placement and safety evidence, not a throughput
graph. The relevant tests cover preferred startup, preferred log replication,
no-op/watermark behavior, and failover/failback when available. In the thesis
body, these tests belong in a compact validation matrix rather than a
line-by-line walkthrough.
